% Tabasco generation metrics -- \input{} from report-draft.tex Ch5.
\begin{table}[tbp]
\centering
\small
\begin{tabular}{l c c c c c c c}
\toprule
\textbf{Model} & \textbf{Steps} & \textbf{Valid.} & \textbf{Conn.} & \textbf{Uniq.} & \textbf{Div.} & \textbf{Nov.} & \textbf{FCD}~$\downarrow$ \\
\midrule
Baseline                     & 151k & 0.980 & 0.998 & 1.000 & 0.886 & 0.966 & \textbf{5.61} \\
\midrule
REPA-CheMeleon               & 73k  & 0.972 & 0.999 & 1.000 & 0.882 & 0.961 & 7.43 \\
REPA-MACE                    & 69k  & 1.000 & 0.995 & 0.996 & 0.884 & 0.977 & 6.81 \\
\bottomrule
\end{tabular}
\caption{\textbf{REPA brings no measurable gain over the baseline.} The quality metrics are saturated, sitting at or near their ceiling for baseline and variants alike. On FCD, the only metric with headroom, both variants are worse. (1{,}000 sampled molecules; higher-is-better except FCD ($\downarrow$). \emph{Steps} is the training step at evaluation; the baseline trained ${\sim}2{\times}$ longer, but the quality metrics plateau early, so the unequal steps do not make the comparison unfair (\S\ref{sec:tabasco-results}). One run per variant.)}
\label{tab:tabasco-gen}
\end{table}
